% !TEX TS-program = pdflatex
% !TEX encoding = UTF-8 Unicode

% Article template :)

% General stuff
\documentclass[11pt]{article} % use larger type; default would be 10pt

\usepackage[utf8]{inputenc} % set input encoding (not needed with XeLaTeX)
\usepackage[french]{babel} % Language setting
\usepackage [T1] {fontenc}
\pdfmapfile{Overpass.map}

% Geometry
\usepackage{geometry}
\geometry{a4paper}
\geometry{margin=2cm}

% Usefull packages
\usepackage{eso-pic} % to add background images to the article
\usepackage{graphicx} % support the \includegraphics command and options
%\usepackage[colorlinks=false]{hyperref}
\usepackage{titlesec} % to change default styles of /sections, /title ...
\usepackage{xcolor}
\usepackage{tabto}
%\usepackage{amssymb}
\usepackage{amssymb}
\usepackage{multirow}
\usepackage{tikz}
\usetikzlibrary{arrows,automata,positioning}
\usepackage{amsmath}
\usepackage{lettrine}
\usepackage{oldgerm}
\usepackage{yfonts}
\usepackage{multicol}
\usepackage{blindtext}
\usepackage{svg}
\usepackage{tabularx}
\usepackage{hyperref}
\hypersetup{
    colorlinks=true,
    linkcolor=blue,
    filecolor=magenta,
    urlcolor=cyan,
    pdftitle={Overleaf Example},
    pdfpagemode=FullScreen,
    }



% Style !
\definecolor{10Black}{cmyk}{0, 0, 0, 0.10} % 10% of black
\definecolor{20Black}{cmyk}{0, 0, 0, 0.20} % 20% of black
\definecolor{80Black}{cmyk}{0, 0, 0, 0.80} % 80% of black
\definecolor{95Black}{cmyk}{0, 0, 0, 0.95} % 60% of black
%\color{10Black}
%\pagecolor{black}

% New commands
\newcommand{\myjump}[1][1]{\mbox{}\\[#1cm]}
\newcommand{\scbf}[1]{\textsc{\textbf{#1}}}
\newcommand{\enluminure}[2]{\lettrine[lines=3]{\small \initfamily #1}{#2}}

\begin{document}
\pagestyle{empty}

\begin{center}
    \textbf{Petit JDR universel}

    Règles de base
\end{center}

\enluminure{P}{eut} importe la quantité de règles qu'un jdr offre, on attendras toujours du MJ qu'il juge les situations non régulées, qu'il improvise et applique la règle du cool. Le petit JDR universel propose peu de règles, il se repose sur les accords qui existent entre joueurs et MJ.


\myjump[0]
\textbf{\huge Joueurs}

\scbf{Arbre de progression}


\noindent
\begin{tabularx}{\textwidth}{cXcX}
\hline
    \scalebox{0.92}{\includegraphics{images/heart.png}} & PV max + 1 &
    \scalebox{0.05}{\includegraphics{images/chestplate.png}} & Defense max + 1\\

    \scalebox{0.03}{\includegraphics{images/cible.png}} & Objectif de Combat (OC) + 1 &
    \scalebox{0.025}{\includegraphics{images/sword.png}} & Augmente le bonus de dégâts\\


    \scalebox{0.03}{\includegraphics{images/esquive.png}} & Esquive + 1, l'esquive augemente de 1 l'OC d'un ennemy vous attaquant&
    \scalebox{0.015}{\includegraphics{images/chance.png}} & Chance + 1, un point de chance permet de relancer un dés. Le MJ décide quand est-ce que les points de chance se reset\\


    \vspace{0.3cm}\scalebox{0.065}{\includegraphics{images/pouch.png}} & \multicolumn{3}{@{}l}{Augmente de 1 le nombre d'armes et objets pouvant être utilisés sans malus en combat}\\

\hline
    \textbf{PV} &Lorsque le personnage atteint 0 PV, il doit recevoir un coup supplémentaire avant de mourrir &
    \textbf{Stress} & Par défaut à 0/6, augmente au fur et à mesure de la partie. C'est le MJ qui remet le stress à 0\\
\hline
\end{tabularx}

\myjump[0.3]
\indent\scbf{Compétences}\newline
À part dans l'arbre de progression, il n'existe pas de compétences pré-définies, c'est au joueur d'imaginer les compétences que son personnage possède. Pour réaliser des actions hors combat, le joueur doit justifier au MJ pourquoi son personnage est capable de réaliser la tache envisagée. Cette justification peut utiliser comme argument l'une des compétences du personnage.
\myjump[0.3]
\begin{tabularx}{\linewidth}{|X}
\emph{\og J'aimerais que mon personnage cherche du cerfeuille dans la forêt. Il devrait en être capable d'en trouver car il a été élève d'un druide et possède la compétence herboristerie 2. \fg}\\
\end{tabularx}
\myjump[0]
\textbf{Ignore PV7/S3:} Ignore l'effet appliqué quand vous atteignez 7 de PV et 3 de stress.\newline
\textbf{Parkour, Discret, Puissant, Intimidation:} Permet de justifier au MJ d'escalader une facade, défoncer une porte, être discret ou intimidant.



\myjump[0]
\indent\scbf{Création de personnage}\newline
Esquive et chance sont à 0, défense max à 4, PV max à 12 et OC à 2. Distribuer 10 points d'expérience dans l'arbre de progression. 6 points de compétences sont disponibles et peuvent être distribués plus tard dans la partie. Choisir une armes de départ à 4 de dégâts (3 si arme distance) et de 2 armures qui s'additionnent à votre défense max-2. Le MJ à le dernier mot vis à vis de votre équipement

\myjump[0]
\indent\scbf{Combat}\newline
En combat, votre personnage peut réaliser 2 actions dont une seule qui lance ses 2d6 de dégâts. Par exemple, attaquer puis pousser une bibliothèque sur un ennemi, les 2 actions font des dégâts mais seulement la première lance les 2d6, la seconde c'est au MJ d'interpreter.
Lorsque votre personnage veut attaquer, il n'y a qu'un seul jet de dés à faire pour connaître les dégâts qu'il infligera. Vous lancez 2d6 explosifs que vous comparez à votre Objectif de Combat (OC). Si le résultat est inférieur ou égal à votre OC, vous additionnez au jet de dés vos dégâts bonus et les dégâts de votre arme.

$
\mbox{Dégâts infligés} =
\left\{
    \begin{array}{ll}
        (2d6 \le \mbox{OC}) + \mbox{dégâts magiques} + \mbox{degâts bonus} + \mbox{dégats arme}\\
(2d6 > \mbox{OC}) + \mbox{dégâts magiques}
    \end{array}
\right.
$

\noindent Lorsque votre personnage se fait attaquer, vous lancez un d6. Si le résultat est supérieur à votre esquive, le MJ vous indique les dégâts \og haut\fg de l'ennemi qui vous attaque, sinon, ce sont les dégâts \og bas\fg.










\newpage
\textbf{\huge Maître du Jeu}

\enluminure{L}{'improvisation} ne peut pas être utilisé pour solidifier du gameplay. Si la magie fait partie de votre univers mais qu'il n'y a pas de règles qui explique son fonctionement; alors cette magie est un outil scénaristique pour le MJ plutôt qu'un élément de gameplay pour les joueurs.

\myjump[0]
\indent\scbf{Modules}\newline
\noindent Pour ajouter des éléments de gameplay au set de règles de ce jeu de rôle, vous pouvez ajourter des modules que vous retrouverez sur la page suivante : \href{https://github.com/kalharko/petit-jdr-universel}{github.com/kalharko/petit-jdr-universel}.



\myjump[0]
\indent\scbf{Compétences}\newline
Le cadre de \emph{compétence} de la feuille de personnage recueille tous les domaines où le personnage possède une maîtrise ou un savoir. Pour réaliser des actions hors combat, les joueurs vont justifier pourquoi leur personnage est capable de réaliser la tache envisagée. Vous devez juger leur argumentaire en prenant en compte les compétences du personnage, sa condition et l'histoire. Les domaines techniques sont très faciles à juger, quelqu'un qui ne s'est jamais entraîné au crochetage, ne peut pas crocheter. Pour les situations plus nuancés le MJ doit prendre en compte l'histoire et choisir le jugement qui va ajouter le plus de drame et de rebondissement à l'histoire.

\myjump[0]
Il est conseillé d'offrir des points de temps à vos joueurs, qui représentent le temps libre dont disposent les personnages entre 2 aventures. Ils peuvent les dépenser pour crafter un objet, améliorer ou apprendre une compétence ou un coup spécial.

\myjump[0]
\indent\scbf{Coups spéciaux}\newline
Un coups spécial est une compétence utilisable une fois par combat qui inflige des dégâts ou un effet à un ennemi. C'est au joueurs de les inventer pour leurs personnages, et au MJ de réguler leurs effets. Le dernier mot revient toujours au MJ, il décide même si l'effet est appliqué à son boss ou non.

\myjump[0]
\indent\scbf{Équipement}\newline
Les armes à distance se jouent comme les autres armes, mais elles infligent légèrement moins de dégâts. Les boucliers possèdent une faible attaque et un bonus de défense, pour pouvoir en profiter il faut avoir utilisé une action à son tour pour lever son bouclier. Le bonus de défense d'un bouclier permet de dépasser sa défense max.


\myjump[0]
\indent\scbf{Combat}\newline
Lorsque vous préparez un combat, les ennemis sont résumés par les informations suivantes :\newline
\begin{tabularx}{\linewidth}{c|X}
\hline
    \scbf{Nom/Titre} & \emph{Guerrier Insecte}\\\hline
    \multicolumn{2}{@{}l}{
    \begin{tabularx}{\linewidth}{cXcXcX}
        \scbf{Degats bas} & \emph{8} & \scbf{Degats haut} & \emph{12} & \scbf{PV} & \emph{20}\\\hline
    \end{tabularx}
    }\\
    Description et & \emph{Un insecte à taille humaine, qui défends sa reine à tout prix.}\\
    Coups spéciaux & \emph{Une fois par combat,  attaque à 14 de dégâts haut avec ses mandibules.}\\
\hline
\end{tabularx}

\myjump[0.35]
Les dégâts haut, bas et PV peuvent être déterminés grâce au tableur d'équilibrage fournit avec ce jeu de rôle. Il ne vous reste plus qu'à imaginer la personnalité de la créature et ses coups spéciaux. Cette méthode suis le paradigme de combat orienté action comme décrit par Matt Colvile : \href{https://youtu.be/y_zl8WWaSyI}{vidéo youtube}.



\end{document}
